\section{Conclusion}
\label{sec:conclusion}

We built a ticket ordering system for a trilingual bank support desk and
evaluated it where it is used, in a queue, rather than where it is convenient to
measure, in a classification table. The Ticket Urgency Score combines a priority
posterior, a sentiment posterior, and a learned per-intent criticality under
multiplicative aging, and attains 97--99\% of the reduction in urgent-ticket
waiting time that a perfect classifier would deliver.

The finding we would put in front of an operations owner is that triage fairness
is a scheduling property. A naive argmax deployment makes customers who write
romanized Sinhala or Tamil wait 4.07 and 6.99 minutes longer for the same urgent
problem, with the gap absent under first-come-first-served, absent under a
gold-label control, and absent on both native-script tracks. Reading the same
model's posterior through our score cuts it to $+0.37$ and $+0.50$.

Two negative results travel with it. Rank correlation selects the worst system
in our comparison, because a queue is consumed from the top and $\tau$ averages
over the whole list. And classifier accuracy does not transfer: a two-point
macro-F\textsubscript{1} advantage leaves the median wait unchanged and produces
a six-hour tail, because the stronger model makes fewer errors but
$8.5\times$ more confident ones.

The corpus, the frozen split, and the model checkpoints are publicly available
under CC BY 4.0 with attribution to BANKING77.
\ifblind\relax\note{Repository URLs withheld for double-blind review; add for
camera-ready.}\fi
